\section{Network Observability}
\label{chap:impl:observability}

The testbed's monitoring system incorporates both the passive and active monitoring methods identified in Section~\ref{chap:proposal:observability}. This section describes the implementation of each monitoring approach, followed by the mechanisms through which the collected metrics are exposed.

\subsection{Passive Monitoring}
\label{chap:impl:observability:passive}

Ideally, monitoring network usage would rely on interfaces already exposed by the core network. However, as the 5GAIner 5G cores do not provide such monitoring APIs, an alternative solution was developed to mitigate this limitation. Traffic metrics are passively collected by a Prometheus exporter deployed on a man-in-the-middle switch positioned on the data path of the 5G system\footnote{\later{update github link https://github.com/Migas77/dissertationInit/tree/main/monitoring}}.

The traffic observed on the switch comprises \acs{5GC} traffic and \acs{BBU} traffic, each forwarded to a separate interface of a dedicated \acs{VM} with restricted access only to privileged users, e.g. \acs{ATNoG} system administrator. To overcome this restriction, the mirrored traffic is instead redirected to a separate, accessible \acs{VM} dedicated to metrics collection. For each monitored interface, packets captured with \texttt{tcpdump}\footnote{https://man7.org/linux/man-pages/man1/tcpdump.1.html} are relayed as a pcap stream, via \texttt{socat}\footnote{https://man7.org/linux/man-pages/man1/socat.1.html}, over a TCP connection to the metrics \acs{VM}. There, a Python script parses the stream with Scapy\footnote{https://scapy.net} and injects the packets into a corresponding \acs{TAP} interface, replicating the traffic as it was received on the original interface. Both the sender and the receiver scripts run as \texttt{systemd}\footnote{https://man7.org/linux/man-pages/man1/systemd.1.html} services, ensuring that both processes are started at boot and automatically restarted, re-establishing the stream after reboots or connection failures. As a result, each  \acs{TAP} interface replicates the traffic of its corresponding switch interface, allowing monitoring tools, i.e. the exporter, to collect network metrics as if they were directly attached to the original switch, and the \acs{5G} data path.

The collected metrics include uplink and downlink traffic volume, in bytes, partitioned by source and destination address (IP and MAC) and the underlying transport protocol. The exporter, implemented in Go, reads the packets arriving at a \acs{TAP} interface with \texttt{gopacket}\footnote{https://pkg.go.dev/github.com/google/gopacket} and \texttt{libpcap}\footnote{https://www.tcpdump.org/manpages/pcap.3pcap.html}, parses their link and network layer headers, and increments byte and packet counters accordingly. One exporter instance is deployed per \acs{TAP} interface, each deployed as a separate Docker container. From these byte counters, throughput metrics are derived by calculating the traffic volume variation over a configurable interval. These metrics are scraped and stored by a Prometheus instance that makes them available for querying and aggregation.

Since the metrics are gathered from traffic captured directly at each point of the 5G network, no additional instrumentation on the UE is required, with the observed values reflecting the actual behavior of the network. Ultimately, this solution can be easily integrated into other testbeds, providing a non-intrusive way of monitoring network traffic and overcome the lack of dedicated interfaces in other open source and commercial core implementations.

\subsection{Active Monitoring}
\label{chap:impl:observability:active}

Complementing this passive monitoring approach, the testbed also supports active monitoring, allowing the generation of synthetic traffic towards provisioned \acs{UE} to assess their connectivity. As identified in Section~\ref{chap:proposal:observability:activeprobing}, this requires a prober component that periodically sends requests towards monitored \acs{UE} and computes performance metrics based on the received responses. These metrics are exposed to and scraped by the same Prometheus instance used for passive monitoring. Although such a component could be developed from scratch, several open source Prometheus exporters already address this problem. Among them, the \texttt{network\_exporter}\footnote{https://github.com/syepes/network\_exporter} and the \texttt{blackbox\_exporter}\footnote{https://github.com/prometheus/blackbox\_exporter}, the latter maintained by the Prometheus project, were identified as the most promising alternatives.

Both solutions fulfill the intended goal, allowing the probing configuration to be changed and updated at runtime, which is required for a testbed where the set of probed targets follows the lifecycle of the provisioned \acs{UE}. However, they operate on fundamentally different principles. The \texttt{blackbox\_exporter} follows an on-demand model issuing a single request each time it is scraped by the Prometheus. In contrast, the \texttt{network\_exporter} is completely decoupled from Prometheus. , probing each target with a burst of echo requests at an interval defined in the prober configuration, with each scrape simply returning the most recent results

This is achieved through a dedicated prober, based on the open source \texttt{network\_exporter}\footnote{https://github.com/syepes/network\_exporter} that sends \acs{ICMP} echo requests and exposes the computed metrics to the same Prometheus instance. For each configured target, the prober probes and collects 

In addition, this prober can also be employed to conduct throughput assessments. Since probing introduces additional load on the network, it is only activated when explicitly requested by the experimenter, being tied to the UE lifecycle to ensure probes are automatically started only after a UE is associated to a slice. This requires integration at the OSS and orchestration level, which is presented later in the section.

\subsection{Metrics Exposure}
\label{chap:impl:observability:exposure}

The measurements collected with both methods are exposed through two mechanisms. Grafana dashboards graphically present network slice and UE metrics, and are automatically provisioned and updated by the OSS as network slice and UE resources change. In parallel, programmatic access through the NEF AnalyticsExposure and CAMARA Connectivity Insights Subscriptions APIs outlined earlier is also provided. To achieve this, the NEF queries the Prometheus HTTP API to derive near-real-time and historic analytics from the stored time series. Similarly, the offering of active monitoring metrics is also achieved through the same workflow.

\plan{
    \begin{enumerate}
        \item (To be continued ...)
        \item Establish
        \item Highlight
        \item Introduce
    \end{enumerate}
}
